% ============ 第1章 §1.5 覆盖定理与Hardy-Littlewood极大函数(原书 pp.37-44 = PDF pp.49-56) ============
\section{覆盖定理与 Hardy--Littlewood 极大函数}\label{sec:1.5}
\index{Hardy--Littlewood 极大函数}
\index{Hardy--Littlewood 中心极大函数}
\index{覆盖定理!Besicovitch 的}
\index{弱 (1,1) 型!Hardy--Littlewood 极大函数的}
\index{Chebyshev 不等式}


为了研究 Hardy--Littlewood 极大函数的基本性质,先来研究覆盖定理.以下,用 $\ell(Q)$ 表示区间 $Q$ 的长度.又,$Q(x, r)$ 表示中心为 $x$,长度为 $r$ 的区间 $(x - r/2, x + r/2)$.

\begin{theorem}[Besicovitch 覆盖定理]\label{thm:1.5.1}
设 $A$ 是 $\mathbb{R}$ 内的有界集,对每个 $x \in A$ 对应一个以 $x$ 为中心的区间 $Q(x)$,且 $\sup_{x \in A}\ell(Q(x)) < \infty$.则可以从区间族 $\{Q(x)\}_{x \in A}$ 中选出至多可数个区间列 $\{Q_{k}\}$,使其满足以下 3 个条件:

\begin{enumerate}[(i)]
\item $A \subset \bigcup_{k}Q_{k}$,
\item $\sum_{k}\chi_{Q_{k}}(x) \leq 4$,即包含 $x$($\in A$)的 $Q_{k}$ 至多 4 个,
\item 区间列 $\{Q_{k}\}$ 可分解为至多 9 组互不相交的子区间列 $\{Q_{k_{1}}\}$,$\{Q_{k_{2}}\}$,$\ldots$,$\{Q_{k_{\ell}}\}$($\ell \leq 9$),即各 $\{Q_{k_{j}}\}$ 互不相交,且 $\{Q_{k_{1}}\} \cup \{Q_{k_{2}}\} \cup \cdots \cup \{Q_{k_{\ell}}\} = \{Q_{k}\}$.
\end{enumerate}
\end{theorem}

\begin{proof}
设 $\alpha_{0} := \sup_{x \in A}\ell(Q(x)) < \infty$.首先选取 $x_{1} \in A$ 使 $\ell(Q(x_{1})) > \frac{1}{2}\sup_{x \in A}\ell(Q(x))$.其次,若 $A \setminus Q(x_{1}) \neq \emptyset$,选取 $x_{2} \in A \setminus Q(x_{1})$ 使 $\ell(Q(x_{2})) > \frac{1}{2}\sup_{x \in A \setminus Q(x_{1})}\ell(Q(x))$.以下归纳地进行:若 $A \setminus \bigcup_{j=1}^{k}Q(x_{j}) \neq \emptyset$,选取 $x_{k+1} \in A \setminus \bigcup_{j=1}^{k}Q(x_{j})$ 使 $\ell(Q(x_{k+1})) > \frac{1}{2}\sup_{x \in A \setminus \bigcup_{j=1}^{k}Q(x_{j})}\ell(Q(x))$.若 $A \setminus \bigcup_{j=1}^{k}Q(x_{j}) = \emptyset$,则在该阶段终止此操作.对如此得到的区间列 $\{Q(x_{k})\}$,令 $Q_{k} = Q(x_{k})$.下面证明这个区间列 $\{Q_{k}\}$ 满足 (i),(ii),(iii).

首先,成立以下 3 条.

(a) 若 $j < k$,则 $x_{k} \notin Q_{j}$,

(b) 若 $j < k$,则 $\ell(Q_{j}) > \frac{1}{2}\ell(Q_{k})$,

(c) $Q\bigl(x_{j}, \ell(Q_{j})/3\bigr) \cap Q\bigl(x_{k}, \ell(Q_{k})/3\bigr) = \emptyset$($j \neq k$).

(a),(b) 由 $Q_{k}$ 的构造方式是显然的.为验证 (c),可设 $j < k$.于是由 (a),$|x_{k} - x_{j}| \geq \ell(Q_{j})/2$;使用 (b),得 $\frac{1}{6}\ell(Q_{j}) + \frac{1}{6}\ell(Q_{k}) < \frac{1}{6}\ell(Q_{j}) + \frac{1}{3}\ell(Q_{j}) = \frac{1}{2}\ell(Q_{j}) \leq |x_{k} - x_{j}|$,即区间 $Q\bigl(x_{j}, \ell(Q_{j})/3\bigr)$ 与 $Q\bigl(x_{k}, \ell(Q_{k})/3\bigr)$ 的中心间的距离大于各自长度之和的一半,故 $Q\bigl(x_{j}, \ell(Q_{j})/3\bigr)$ 与 $Q\bigl(x_{k}, \ell(Q_{k})/3\bigr)$ 没有公共部分.

现在回到 (i) 的证明.若构造区间列 $\{Q(x_{k})\}$ 的操作有限次结束,则 (i) 已证,故设它无限继续下去.此时,$\lim_{j \to \infty}\ell(Q_{j}) = 0$.事实上,若设 $\limsup_{j \to \infty}\ell(Q_{j}) = a > 0$,则存在子列 $\{Q_{j_{k}}\}$ 使得 $\ell(Q_{j_{k}}) > a/2$.令 $Q'_{j_{k}} = Q\bigl(x_{j_{k}}, \frac{1}{3}\ell(Q_{j_{k}})\bigr)$,则
\begin{equation}\tag{1.5.1}
\sum_{k=1}^{\infty}\ell(Q'_{j_{k}}) > \sum_{k=1}^{\infty}\frac{a}{6} = +\infty.
\label{eq:1.5.1}\end{equation}
另一方面,$Q'_{j_{k}} \subset Q_{j_{k}} \subset Q\bigl(x_{1}, \mathrm{diam}(A) + \alpha_{0}\bigr)$,且由 (c),$Q'_{j_{k}}$ 两两不相交,所以
\begin{equation}\tag{1.5.2}
\sum_{k=1}^{\infty}\ell(Q'_{j_{k}}) < \mathrm{diam}(A) + \alpha_{0}.
\label{eq:1.5.2}\end{equation}
式~\eqref{eq:1.5.1} 与式~\eqref{eq:1.5.2} 矛盾.因此 $\lim_{j \to \infty}\ell(Q_{j}) = 0$.

于是,假设存在 $x_{0} \in A \setminus \bigcup_{j=1}^{\infty}Q_{j}$.则由 $\lim_{j \to \infty}\ell(Q_{j}) = 0$,存在 $Q_{j_{0}}$ 使得 $\ell(Q_{j_{0}}) < \frac{1}{2}\ell(Q_{x_{0}})$.由 $x_{0} \in A \setminus \bigcup_{j=1}^{\infty}Q_{j}$,有 $x_{0} \in A \setminus \bigcup_{j=1}^{j_{0}-1}Q_{j}$,故由 $Q_{j_{0}}$ 的定义,
\begin{equation*}
\ell(Q_{j_{0}}) > \frac{1}{2}\sup_{x \in A \setminus \bigcup_{j=1}^{j_{0}-1}Q(x_{j})}\ell(Q(x)) \geq \frac{1}{2}\ell(Q(x_{0})).
\end{equation*}
这与 $\ell(Q_{j_{0}}) < \frac{1}{2}\ell(Q_{x_{0}})$ 矛盾.因此
\begin{equation*}
A \setminus \bigcup_{j=1}^{\infty}Q_{j} = \emptyset \quad \text{即} \quad A \subset \bigcup_{j=1}^{\infty}Q_{j}.
\end{equation*}

(ii) 的证明.由 (c),$x_{i} \neq x_{j}$($i \neq j$).对给定的 $x$,设满足 $x \in Q_{j}$ 且 $x \leq x_{j}$ 的编号 $j$ 至少有 2 个,并设其中最小的编号为 $j_{1}$,第二小的编号为 $j_{2}$.今设 $m > j_{2}$ 且 $x < x_{m}$.由 (a),$x_{j_{2}}, x_{m} \notin Q_{j_{1}}$,$x_{m} \notin Q_{j_{2}}$.又由 (b),
\begin{equation}\tag{1.5.3}
\frac{1}{2}\ell(Q_{m}) < \ell(Q_{j_{1}}).
\label{eq:1.5.3}\end{equation}
进一步,由 $x \in Q_{j_{2}}$,$x \leq x_{j_{1}} < x_{j_{2}}$,$x_{j_{2}} \notin Q_{j_{1}}$,得 $\ell(Q_{j_{2}})/2 > \ell(Q_{j_{1}})/2$.因此,再利用 $x_{j_{2}} \notin Q_{j_{1}}$,$x_{m} \notin Q_{j_{2}}$ 以及 $x \leq x_{j_{1}} < x_{j_{2}} < x_{m}$,得
\begin{equation}\tag{1.5.4}
|x_{m} - x| \geq |x_{m} - x_{j_{2}}| + |x_{j_{2}} - x_{j_{1}}| > \frac{\ell(Q_{j_{2}})}{2} + \frac{\ell(Q_{j_{1}})}{2} > \ell(Q_{j_{1}}).
\label{eq:1.5.4}\end{equation}
由 \eqref{eq:1.5.3},\eqref{eq:1.5.4} 可知 $x \notin Q_{m}$.因此,中心在 $x$ 右侧且包含 $x$ 的 $Q_{j}$ 至多 2 个.同样地,中心在 $x$ 左侧且包含 $x$ 的 $Q_{j}$ 也至多 2 个.所以,包含 $x$ 的 $Q_{j}$ 至多 4 个.

(iii) 的证明.设 $j \in \mathbb{N}$.若 $m < j$ 且 $Q_{m} \cap Q_{j} \neq \emptyset$,则 $Q_{m}$ 要么包含 $Q_{j}$ 的端点,要么不包含端点.若不包含端点,则 $Q_{m} \subset Q_{j}$,且由 (b),$\ell(Q_{m}) > \frac{1}{2}\ell(Q_{j})$,故 $x_{j} \in Q_{m}$.但这与 (a) 矛盾.因此 $Q_{m}$ 包含 $Q_{j}$ 的两个端点之一.于是,满足 $m < j$ 且 $Q_{m} \cap Q_{j} \neq \emptyset$ 的编号 $m$ 由 (ii) 至多为 $4 \times 2 = 8$ 个.于是,把 (i) 的证明中构造的区间列 $\{Q(x_{k})\}$ 按如下方式分成 9 个类 $I_{1}, I_{2}, \ldots, I_{9}$.首先,把 $Q_{1}$ 放入类 $I_{1}$,$Q_{2}$ 放入类 $I_{2}$,$\ldots$,$Q_{9}$ 放入类 $I_{9}$.与 $Q_{10}$ 相交的 $Q_{j}$($j < 10$)至多 8 个,故至少有一个与 $Q_{10}$ 不相交的 $Q_{\ell}$($\ell \leq 9$).把 $Q_{10}$ 放入该编号的类 $I_{\ell}$.其次,与 $Q_{11}$ 相交的 $Q_{j}$($j < 11$)至多 8 个,故存在编号 $\ell$($\ell \leq 9$)使得
\begin{equation*}
Q_{11} \cap \bigcup_{Q_{j} \in I_{\ell},\, j < 11}Q_{j} = \emptyset.
\end{equation*}
把 $Q_{11}$ 放入该编号的类 $I_{\ell}$.继续这一操作,则每个 $Q_{i}$ 都落入 $I_{1}, I_{2}, \ldots, I_{9}$ 中的某一类,且各同类 $I_{\ell}$ 中的 $Q_{j}$ 两两不相交.即,这种分类正是 (iii) 所求的分类.
\end{proof}

利用这个覆盖定理的 (i),(ii),可证明称为 Hardy--Littlewood 极大函数弱 $(1,1)$ 型的定理.

\begin{theorem}\label{thm:1.5.2}
存在常数 $C$,使得对任意 $f \in L^{1}(\mathbb{R})$,
\begin{equation*}
|\{x \in \mathbb{R};\ (M_{1}f)(x) > \lambda\}| \leq C\frac{\| f \|_{L^{1}}}{\lambda} \quad (\lambda > 0).
\end{equation*}
\end{theorem}

\begin{proof}
首先,对中心极大函数 $\tilde{M}_{1}$ 证明
\begin{equation*}
|\{x \in \mathbb{R};\ (\tilde{M}_{1}f)(x) > \lambda\}| \leq 4\frac{\| f \|_{L^{1}}}{\lambda} \quad (\lambda > 0).
\end{equation*}
则由引理~\ref{lem:1.4.4},$M_{1}f(x) \leq 2\tilde{M}_{1}f(x)$,定理即得证.

现在,设 $\lambda > 0$,任意固定 $B > 0$.定义集合 $A_{B}$ 为
\begin{equation*}
A_{B} := (-B, B) \cap \{x \in \mathbb{R};\ (\tilde{M}_{1}f)(x) > \lambda\}.
\end{equation*}
对每个 $x \in A_{B}$,由 $(\tilde{M}_{1}f)(x)$ 的定义,存在 $r_{x} > 0$ 使得
\begin{equation*}
\frac{1}{|Q(x, r_{x})|}\int_{Q(x, r_{x})}|f(y)|\dd y > \lambda,
\end{equation*}
故对应一个区间 $Q(x, r_{x})$.于是,由 Besicovitch 覆盖定理,可从区间族 $\{Q(x, r_{x})\}_{x \in A_{B}}$ 中选出满足
(i) $A_{B} \subset \bigcup_{k}Q_{k}$,
(ii) $\sum_{k}\chi_{Q_{k}}(x) \leq 4$
的区间列 $\{Q_{k}\}$.于是,由 (i),
\begin{align*}
|A_{B}| &\leq \sum_{k}|Q_{k}| \leq \sum_{k}\frac{1}{\lambda}\int_{Q_{k}}|f(y)|\dd y = \frac{1}{\lambda}\int_{\mathbb{R}}\sum_{k}\chi_{Q_{k}}(y)|f(y)|\dd y \\
&\leq \frac{4}{\lambda}\int_{\mathbb{R}}|f(y)|\dd y.
\end{align*}
因此,令 $B \to \infty$,得
\begin{equation*}
|\{x \in \mathbb{R};\ (\tilde{M}_{1}f)(x) > \lambda\}| \leq 4\frac{\| f \|_{L^{1}}}{\lambda},
\end{equation*}
定理证毕.
\end{proof}

\begin{note}\label{note:1.5.3}
令 $f(x) = \chi_{[-1,1]}$,则
\begin{equation*}
M_{1}f(x) = \left\{ \begin{array}{l l}
1 & (|x| \leq 1), \\
\frac{2}{|x| + 1} & (|x| > 1)
\end{array} \right.
\end{equation*}
所以对 $0 < \lambda < 1$,
\begin{equation*}
|\{x \in \mathbb{R};\ (M_{1}f)(x) > \lambda\}| = \frac{4}{\lambda} - 2.
\end{equation*}
\end{note}

\begin{note}\label{note:1.5.4}
定理~\ref{thm:1.5.2} 也可以看作对 Chebyshev 不等式
\begin{equation*}
|\{x \in \mathbb{R};\ |f(x)| > \lambda\}| \leq \frac{\| f \|_{L^{1}}}{\lambda} \quad (\lambda > 0)
\end{equation*}
的精细化.
\end{note}

由定理~\ref{thm:1.5.2} 可导出 Hardy--Littlewood 极大函数的 $L^{p}$ 有界性.

\begin{theorem}\label{thm:1.5.5}
设 $1 < p \leq \infty$.则存在只与 $p$ 有关的正常数 $C_{p}$,使得
\begin{equation*}
\| M_{1}f \|_{L^{p}} \leq C_{p}\| f \|_{L^{p}} \quad (f \in L^{p}(\mathbb{R})).
\end{equation*}
\end{theorem}

\begin{proof}
(i) $p = \infty$ 的情形.若 $f \in L^{\infty}(\mathbb{R})$,则
\begin{equation*}
\frac{1}{|Q|}\int_{Q}|f(y)|\dd y \leq \| f \|_{L^{\infty}}\frac{1}{|Q|}\int_{Q}\dd y = \| f \|_{L^{\infty}}
\end{equation*}
所以对任意 $x \in \mathbb{R}$,
\begin{equation*}
M_{1}f(x) \leq \| f \|_{L^{\infty}}.
\end{equation*}
因此
\begin{equation*}
\| M_{1}f \|_{L^{\infty}} \leq \| f \|_{L^{\infty}} \quad (f \in L^{\infty}(\mathbb{R})),
\end{equation*}
故 $p = \infty$ 时定理成立.

(ii) $1 < p < \infty$ 的情形.设 $f \in L^{p}(\mathbb{R})$.对 $\lambda > 0$,令 $E_{\lambda} = \{x \in \mathbb{R};\ M_{1}f(x) > \lambda\}$,$\sigma(\lambda) = |E_{\lambda}|$,则
\begin{equation*}
\int_{\mathbb{R}}|M_{1}f(x)|^{p}\dd x = -\int_{0}^{\infty}\lambda^{p}\dd \sigma(\lambda) = p\int_{0}^{\infty}\lambda^{p-1}\sigma(\lambda)\dd \lambda.
\end{equation*}
进一步,令
\begin{equation*}
f_{1}(x) = \left\{ \begin{array}{l l}
f(x) & (|f(x)| > \lambda/2) \\
0 & (|f(x)| \leq \lambda/2),
\end{array} \right. \quad f_{2}(x) = f(x) - f_{1}(x)
\end{equation*}
则 $|f(x)| \leq |f_{1}(x)| + \frac{\lambda}{2}$,故 $M_{1}f(x) \leq M_{1}f_{1}(x) + \frac{\lambda}{2}$.因此
\begin{equation*}
E_{\lambda} \subset \left\{x \in \mathbb{R};\ M_{1}f_{1}(x) > \frac{\lambda}{2}\right\}.
\end{equation*}
于是,$\sigma(\lambda) \leq |\{x \in \mathbb{R};\ M_{1}f_{1}(x) > \lambda/2\}|$,故使用定理~\ref{thm:1.5.2},
\begin{equation*}
\sigma(\lambda) \leq \frac{16}{\lambda}\int_{\mathbb{R}}|f_{1}(y)|\dd y = \frac{16}{\lambda}\int_{\{y;\ |f(y)| > \lambda/2\}}|f(y)|\dd y.
\end{equation*}
因此
\begin{align*}
\int_{\mathbb{R}}|M_{1}f(x)|^{p}\dd x & = p\int_{0}^{\infty}\lambda^{p-1}\sigma(\lambda)\dd \lambda \\
& \leq 16p\int_{0}^{\infty}\lambda^{p-1}\left(\frac{1}{\lambda}\int_{\{y;\ |f(y)| > \lambda/2\}}|f(y)|\dd y\right)\dd \lambda \\
& = 16p\int_{0}^{\infty}\lambda^{p-2}\left(\int_{\{y;\ |f(y)| > \lambda/2\}}|f(y)|\dd y\right)\dd \lambda \\
& = 16p\int_{\mathbb{R}}|f(y)|\left(\int_{0}^{2|f(y)|}\lambda^{p-2}\dd \lambda\right)\dd y \quad \text{(Fubini)} \\
& = \frac{16p}{p - 1}\int_{\mathbb{R}}2^{p-1}|f(y)|^{p}\dd y.
\end{align*}
由此,
\begin{equation*}
\| M_{1}f \|_{L^{p}} \leq 2\left(\frac{8p}{p - 1}\right)^{1/p}\| f \|_{L^{p}}.
\end{equation*}
\end{proof}

由这个定理可得出 $p$ 次 Hardy--Littlewood 极大函数的 $L^{r}$ 有界性.

\begin{corollary}\label{cor:1.5.6}
设 $1 \leq p < r < \infty$.则存在只与 $p$,$r$ 有关的正常数 $C_{p,r}$,使得
\begin{equation*}
\| M_{p}f \|_{L^{r}} \leq C_{p,r}\| f \|_{L^{r}} \quad (f \in L^{r}(\mathbb{R})).
\end{equation*}
\end{corollary}

\begin{proof}
由 $M_{p}$ 的定义与定理~\ref{thm:1.5.5},
\begin{align*}
\left(\int_{\mathbb{R}}M_{p}f(x)^{r}\dd x\right)^{1/r} &= \left[\left(\int_{\mathbb{R}}M_{1}(|f|^{p})(x)^{r/p}\dd x\right)^{p/r}\right]^{1/p} \\
&\leq \left[C_{r/p}\left(\int_{\mathbb{R}}\left(|f(x)|^{p}\right)^{r/p}\dd x\right)^{p/r}\right]^{1/p} \\
&= (C_{r/p})^{1/p}\| f \|_{L^{r}}.
\end{align*}
\end{proof}
